% Generated draft: requires mathematical and source review.
\documentclass[11pt]{article}
\usepackage[margin=25mm]{geometry}
\usepackage{fontspec}
\setmainfont{Latin Modern Roman}
\usepackage{amsmath,amssymb,mathtools,unicode-math}
\setmathfont{Latin Modern Math}
\usepackage{enumitem,needspace,xurl,hyperref}
\setlength{\parindent}{0pt}
\setlength{\parskip}{5pt}
\setlength{\emergencystretch}{3em}
\newcommand{\N}{\mathbb{N}}
\newcommand{\Z}{\mathbb{Z}}
\newcommand{\Q}{\mathbb{Q}}
\newcommand{\R}{\mathbb{R}}
\newcommand{\C}{\mathbb{C}}
\newcommand{\F}{\mathbb{F}}
\newcommand{\A}{\mathbb{A}}
\newcommand{\PP}{\mathbb{P}}
\newcommand{\E}{\mathbb{E}}
\newcommand{\eps}{\varepsilon}
\newcommand{\dd}{\,\mathrm{d}}
\begin{document}
\section*{UnsolvedMath ID: 30001481}
\textbf{Source problem code:} catalogue entry
\subsection{Definitions and mathematical statement}
For a standard mathematical framework in which to read the question, let \((X,d)\) be a complete separable metric space and fix \(x_0\in X\). Write
\[
\mathcal P_2(X)=\left\{\mu:\mu\text{ is a Borel probability measure on }X,\quad \int_X d(x,x_0)^2\,\mathrm d\mu(x)<\infty\right\}.
\]
For \(\mu,\nu\in\mathcal P_2(X)\), let \(\Pi(\mu,\nu)\) denote the Borel probability measures on \(X\times X\) with first marginal \(\mu\) and second marginal \(\nu\). The \(2\)-Wasserstein distance is
\[
W_2(\mu,\nu)=\left(\inf_{\pi\in\Pi(\mu,\nu)}\int_{X\times X}d(x,y)^2\,\mathrm d\pi(x,y)\right)^{1/2}.
\]
If \(N\geq1\) particles have positions \(X_1(t),\ldots,X_N(t)\in X\) at time \(t\), their empirical probability measure is
\[
\mu_t^N=\frac1N\sum_{i=1}^N\delta_{X_i(t)},
\]
where \(\delta_x\) is the unit point mass at \(x\). For finite particle positions this measure belongs to \(\mathcal P_2(X)\). These definitions provide a conventional framework; the supplied source does not prescribe the state space or metric. In the usual Euclidean realization of its example, the particle positions are independent Brownian motions, but the dimension, initial conditions, and boundary conditions are unspecified.

\textbf{Source question.} What physical meaning can be assigned to the \(2\)-Wasserstein distance for particle systems, in particular for empirical measures of many independent Brownian motions?

The supplied record poses an interpretive research question, rather than a quantified conjecture. It does not identify the two measures to be compared, a physical observable or operational meaning to be recovered, or a limiting regime for the number of particles or the observation times. Thus no unique theorem-level target or criterion for a satisfactory physical interpretation can be extracted without additional specifications.
\subsection{Short English statement}
The question asks how to understand the \(2\)-Wasserstein distance physically when probability measures describe particle configurations, especially configurations of many independent Brownian particles. The source leaves the comparison measures and the intended physical interpretation unspecified, so this draft preserves the question as an interpretive research problem.
\subsection{Sources}
\begin{itemize}
\item Local UnsolvedMath registry: supplied record with permanent ID \(30001481\), title \emph{Physical Interpretation of Wasserstein Distance in Particle Systems}, source problem code \texttt{catalogue entry}, in the \emph{Oberwolfach Reports Open Problems} collection. The recorded proposed year is 2010; no individual proposer is supplied.
\item Supplied primary report citation: \emph{Phase Transitions} (2010), Oberwolfach Reports, DOI 10.4171/owr/2010/24, pp.~1431--1484. \url{https://doi.org/10.4171/owr/2010/24}. The collection describes Oberwolfach Reports as documenting mathematical workshops and their open-problem discussions. No page specific to this question is supplied.
\item Supplied literature source, labeled ``OWR 2010/24'': \url{https://doi.org/10.4171/owr/2010/24}.
\item Supplied literature source, labeled ``Wasserstein methods in kinetics'': \url{https://link.springer.com/article/10.1007/s00205-021-01705-9}. No further bibliographic details or precise supporting result are supplied.
\item Literature assessment dated exactly \texttt{2026-08-21}: the supplied curation labels the problem open and states that Wasserstein distance has rigorous roles in particle-system limits and transport estimates, while the broad physical-meaning question is not a single resolved theorem. This is supplied dated curation derived from the OWR collection, not an independently verified literature assessment; its time-sensitive status and cited claims require independent checking.
\item The registry labels the statement ``exact'' and reports that its mathematical content was checked against the official Oberwolfach report, with notation, context, and question form normalized without changing the source-level claim. This verification description is supplied provenance and has not been independently verified for this draft.
\end{itemize}
\end{document}
